\subsection{实测度。正测度}

设 X 是一个局部紧空间。实向量空间 \(\mathcal{K}(\mathrm{X};\mathbf{R})\) 是复向量空间 \(\mathcal{K}(\mathrm{X};\mathbf{C})\) 的底实向量空间的一个子空间；此外，映射 \((f_{1},f_{2})\mapsto f_{1}+if_{2}\) 是积拓扑向量空间 \(\mathcal{K}(\mathrm{X};\mathbf{R})\times\mathcal{K}(\mathrm{X};\mathbf{R})\) 到实拓扑向量空间 \(\mathcal{K}(\mathrm{X};\mathbf{C})\) 上的一个同构（No. 1, 命题 1）。

对每个（复）测度 \(\mu\in\mathcal{M}(\mathrm{X};\mathbf{C})\) ，限制 \(\mu_{0}\) （\(\mu\) 在 \(\mathcal{K}(\mathrm{X};\mathbf{R})\) 上的）是 \(\mathcal{K}(\mathrm{X};\mathbf{R})\) 到 C 中的一个连续 R-线性映射；此外，这一限制确定了 \(\mu\) ，因为若 \(f=f_{1}+if_{2}\) ，其中 \(f_{1},f_{2}\) 属于 \(\mathcal{K}(\mathrm{X};\mathbf{R})\) ，则 \(\mu(f)=\mu_{0}(f_{1})+i\mu_{0}(f_{2})\) 。反之，设 \(\mu_{0}\) 是 \(\mathcal{K}(\mathrm{X};\mathbf{R})\) 到 C 中的一个连续 R-线性映射；显然，映射

\[
f _ {1} + i f _ {2} \mapsto \mu_ {0} (f _ {1}) + i \mu_ {0} (f _ {2})
\]

是 X 上的一个（复）测度。这样，X 上的每个测度可以与它在 \(\mathcal{K}(\mathrm{X};\mathbf{R})\) 上的限制等同起来。

设 \(\mu\) 是 X 上的一个测度。称 \(\mu\) 的共轭测度是这样的测度 \(\overline{\mu}\) ：它由 \(\overline{\mu}(f) = \overline{\mu(\overline{f})}\) （对每个函数 \(f \in \mathcal{K}(\mathrm{X};\mathbf{C})\) ）定义；因为显然 \(\overline{\mu}\) 是一个 C-线性形式，且它在 \(\mathcal{K}(\mathrm{X};\mathbf{C})\) 上连续；显然 \(\overline{\overline{\mu}} = \mu\) ，且对两个测度 \(\mu, \nu\) 及两个标量 \(\alpha, \beta\) （\(\mathbf{C}\) 中的），

\[
\overline {{(\alpha \mu + \beta \nu)}} = \overline {{\alpha}} \cdot \overline {{\mu}} + \overline {{\beta}} \cdot \overline {{\nu}}.
\]

更一般地，对每个函数 \(g \in \mathcal{C}(\mathrm{X};\mathbf{C})\) 及每个测度 \(\mu\) （\(\mathbf{X}\) 上的），

\[
\overline {{{g \cdot \mu}}} = \overline {{{g}}} \cdot \overline {{{\mu}}},\tag{7}
\]

这由定义（No. 4）立即可得。

X 上的测度 \(\mu\) 称为实测度，如果 \(\overline{\mu} = \mu\) ；由前所述，这等价于说对每个函数 \(f \in \mathcal{K}(X; \mathbf{R})\) ，\(\mu(f)\) 是实数。若把实测度与它在 \(\mathcal{K}(X; \mathbf{R})\) 上的限制等同起来，则可以说 X 上的实测度所成的集是实局部凸空间 \(\mathcal{K}(X; \mathbf{R})\) 的对偶；它是一个实向量空间，记作 \(\mathcal{M}(X; \mathbf{R})\) （如不致混淆，有时也记作 \(\mathcal{M}(X)\) ）。R 上的 Lebesgue 测度是实测度，Dirac 测度 \(\varepsilon_{a}\) （对每个点 \(a \in X\) ）亦然。若 \(g \in \mathcal{C}(X; \mathbf{R})\) 且 \(\mu\) 是实测度，则由 (7)，\(g \cdot \mu\) 也是实测度。

设 \(\mu\) 是 X 上的一个（复）测度。由前述定义可知，测度 \(\mu_{1} = (\mu +\overline{\mu}) / 2\) 与 \(\mu_{2} = (\mu -\overline{\mu}) / 2i\) 都是实的；它们分别称为 \(\mu\) 的实部与虚部，分别记作 \(\mathcal{R}\mu\) 与 \(\mathcal{I}\mu\) ；这些测度还可由下述事实刻画：对每个函数 \(f\in \mathcal{K}(\mathbf{X};\mathbf{R})\) ，

\[
\mu_ {1} (f) = \mathcal {R} (\mu (f)), \quad \mu_ {2} (f) = \mathcal {I} (\mu (f)).
\]

显然

\[
\mu = \mu_ {1} + i \mu_ {2}, \quad \overline {{{{\mu}}}} = \mu_ {1} - i \mu_ {2}.
\]

空间 \(\mathcal{K}(\mathrm{X};\mathbf{R})\) ——\(\mathrm{X}\) 上具有紧支集的连续实值函数所成的——显然对序关系 \(f\leqslant g\) 而言显然是一个 Riesz 空间。我们说实测度 \(\mu\) （\(\mathrm{X}\) 上的）是正的，如果 \(\mu (f)\geqslant 0\) 对每个函数 \(f\geqslant 0\) （属于 \(\mathcal{K}(\mathrm{X};\mathbf{R})\) 的）成立；于是它是 Riesz 空间 \(\mathcal{K}(\mathrm{X};\mathbf{R})\) 上的一个正线性形式（Ch. II, §2, No. 1, 定义 1）。反之：

定理 1. --- Riesz 空间 \(\mathcal{K}(\mathrm{X};\mathbf{R})\) 上的每个正线性形式是 \(\mathrm{X}\) 上的一个（正）实测度。

因为，设 \(\mu\) 是 \(\mathcal{K}(\mathrm{X};\mathbf{R})\) 上的一个正线性形式，且设 \(\mathrm{K}\) 是 \(\mathrm{X}\) 的一个紧子集。存在一个具有紧支集的连续映射 \(f_0\) （\(\mathrm{X}\) 到 {[}0, 1{]} 中的），使得 \(f_0(x) = 1\) 在 \(\mathrm{K}\) 上成立 （No. 2, 引理 1）。于是对每个函数 \(g \in \mathcal{K}(\mathrm{X},\mathrm{K};\mathbf{R})\) ，有 \(-\| g \| f_0 \leqslant g \leqslant \| g \| f_0\) ，从而 \(|\mu(g)| \leqslant \| g \| \cdot \mu(f_0)\) ，这就证明了定理。

我们用 \(\mathcal{M}_{+}(X)\) 表示 X 上正测度所成的带顶点的凸锥（换言之，即 Riesz 空间 \(\mathcal{K}(X;\mathbf{R})\) 上正线性形式所成的锥）。

定理 2. --- 局部紧空间 X 上的每个实测度是两个正测度之差。

鉴于定理 1 与 Ch. II, §2, No. 2, Th. 1，只须证明 X 上的实测度 \(\mu\) 是 Riesz 空间 \(\mathcal{K}(\mathrm{X};\mathbf{R})\) 上的一个相对有界线性形式。设 \(f\) 是 X 上的一个连续函数 \(\geqslant 0\) ，具有紧支集 K；关系 \(0 \leqslant g \leqslant f\) （\(\mathcal{K}(\mathrm{X};\mathbf{R})\) 中的）蕴涵 \(\|g\| \leqslant \|f\|\) 且 \(g\) 的支集含于 K。根据假设，存在一个数 \(\mathrm{M}_{\mathrm{K}} \geqslant 0\) ，使得 \(|\mu(h)| \leqslant \mathrm{M}_{\mathrm{K}} \cdot \|h\|\) 对每个函数 \(h \in \mathcal{K}(\mathrm{X},\mathrm{K};\mathbf{R})\) 成立；因此 \(|\mu(g)| \leqslant \mathrm{M}_{\mathrm{K}} \cdot \|g\| \leqslant \mathrm{M}_{\mathrm{K}} \cdot \|f\|\) ，这就证明了定理。

这样，X 上实测度所成的空间 \(\mathcal{M}(\mathrm{X};\mathbf{R})\) 与 Riesz 空间 \(\mathcal{K}(\mathrm{X};\mathbf{R})\) 上相对有界线性形式所成的空间恒同；我们回想，在 \(\mathcal{M}(\mathrm{X};\mathbf{R})\) 中，序关系 \(\mu\leqslant\nu\) 表示 \(\nu-\mu\) 是一个正测度，也表示 \(\mu(f)\leqslant\nu(f)\) 对每个函数 \(f\in\mathcal{K}_{+}(\mathrm{X})\) 成立。

定理 3. --- 局部紧空间 X 上实测度所成的空间 \(\mathcal{M}(\mathrm{X};\mathbf{R})\) 是一个完全格序空间。

这由 Ch. II, §2, No. 2, Th. 1 得出。

按照 Ch. II, §1 的记号，我们对每个实测度 \(\mu\) （\(X\) 上的）定义

\[
\mu^ {+} = \sup (\mu , 0), \quad \mu^ {-} = \sup (- \mu , 0), \quad | \mu | = \sup (\mu , - \mu);
\]

于是 \(\mu = \mu^{+} - \mu^{-}\) ，\(|\mu| = \mu^{+} + \mu^{-}\) 且 \(\inf (\mu^{+},\mu^{-}) = 0\) 。此外，对每个函数 \(f\in \mathcal{K}_{+}(\mathrm{X})\) ，

\[
\int f d \mu^ {+} = \sup _ {0 \leqslant g \leqslant f, g \in \mathcal {K} (\mathrm{X})} \int g d \mu\tag{8}
\]

以及

\[
\int f d | \mu | = \sup _ {| g | \leqslant f, g \in \mathcal {K} (\mathrm{X})} \int g d \mu ,\tag{9}
\]

由此特别得到，

\[
\left| \int f d \mu \right| \leqslant \int | f | d | \mu |\tag{10}
\]

对每个函数 \(f \in \mathcal{K}(\mathbf{X};\mathbf{R})\) 成立。

若 \(f \in \mathcal{K}(\mathrm{X};\mathbf{C})\) ，这一不等式也为真；因为用一个绝对值为 1 的复数去乘 f（这不改变不等式的任何一边），可以假设 \(\int f d\mu \geqslant 0\) 。于是

\[
\left| \int f d \mu \right| = \int f d \mu = \int (\mathcal {R} f) d \mu \leqslant \int | \mathcal {R} f | d | \mu | \leqslant \int | f | d | \mu |.
\]

